\section{Experiments}

\subsection{Evaluation}
In experiments we evaluate based on perplexity (per-wordpiece), a common language
modeling metric, and ROUGE-L F1 (version \texttt{ROUGE-1.5.5}), a common metric used in comparing
candidate and reference summaries. Note the F1 flavor of ROUGE is more 
appropriate in this setting as we do not explicitly constrain the output length in abstractive models;
it is the harmonic
mean of ROUGE-Recall (which favors long summaries) and ROUGE-Precision (which favors short
summaries).

Although optimizing ROUGE directly has been shown to not always yield the best summaries as evaluated by
human judgment \citep{paulus2017deep}, we found that for our task optimizing for perplexity
correlates with increased ROUGE and human judgment. We suspect that the relatively uniform style
of Wikipedia articles makes ROUGE more appropriate here than in general abstractive summarization tasks.

\subsection{Model Training Details and Decoding}
For all abstractive model training, we use the open-source
\texttt{tensor2tensor}\footnote{\url{https://github.com/tensorflow/tensor2tensor}}
library. 

The seq2seq baseline had a hidden size of 128 with 2 layers (we use the hyper-parameter
set defined in the library as \texttt{lstm\_attention}).

For the Transformer encoder-decoder (T-ED), we use the hyper-parameter
set  \texttt{transfomer\_base\_v1}
and train for 1 million steps.
Models exhibited very little over-fitting and did not require early-stopping.
The Transformer Decoder (T-D) was identical to the decoder part of T-ED.
The T-DMCA model is similar to T-D, but with the enhancements described in
section \ref{t-dmca}.

Unless otherwise stated, during decoding we use a beam search of size 4 and length penalty $\alpha=0.6$
\citep{wu2016google} and decode until an end-of-sequence token is reached.

\subsection{Results and Discussion}
\label{discussion}
There are four main dimensions we vary in experiments in generating Wikipedia lead sections:
\begin{enumerate}
    \item Extractive method: \textit{SumBasic}, \textit{TextRank}, \textit{tf-idf}, \textit{identity}, \textit{cheating extractor}
	\item Input corpus: \textit{citations}, \textit{search results}, \textit{combined}
	\item Abstractive model input length, $L$: We try values between 100 and 11000.
	\item Abstractive model architecture: \textit{seq2seq-att}, \textit{T-ED}, \textit{T-D}, \textit{T-DMCA}
\end{enumerate}

\begin{figure}[h]
\begin{center}
\fbox{%
    \includegraphics[width=0.6\textwidth]{extractive_barplot.png}
}%
\caption{ROUGE-L F1 for various extractive methods. The abstractive model contribution is shown
for the best combined \textit{tf-idf}-T-DMCA model.}
\label{extractive_barplot}
\end{center}
\end{figure}
\textbf{Extractive-only is not enough:}
We investigate performance of extractive methods without the abstractive model by looking at the ROUGE-L F1 scores after running \textit{tf-idf}, \textit{SumBasic}, and \textit{TextRank} in
Figure \ref{extractive_barplot}, without any abstractive model.
In the case of TextRank and SumBasic we matched
the output length to the target length and observe the extractive methods perform roughly
in-line with each other in terms of ROUGE-L F1. Our best abstractive model more than doubled
this metric. Further, this model yields large improvements in perceived linguistic quality
(elaborated below).


\textbf{Extractive method:}
From Table \ref{ext-input-perf} we observe that smart extraction is critical
for final abstractive performance.
There is a significant gap between doing nothing, \textit{identity},
and extractive summarization, \textit{tf-idf}. Further, there is a
significant gap between \textit{tf-idf} and the \textit{cheating} extractor, suggesting
future work in improving the extraction step could result in significant improvements.
One possibility is to train a supervised model to predict relevance (Eq. \ref{cheating-similarity}),
which we leave as future work.
For subsequent experiments we fix the extractive method to \textit{tf-idf}.

% Results pulled from
% https://docs.google.com/document/d/1wTI2zzkr39KOclO8h8sFAKtxQeqiAwe0QGaXRvottVc/edit
\begin{table}[t]
\caption{Comparison of extractive method and corpus with $L=500$, and
	the Transformer E-D model}
\label{ext-input-perf}
\begin{center}
\begin{tabular}{llll}
\multicolumn{1}{l}{\bf Extractor}
	&\multicolumn{1}{l}{\bf Corpus}
	&\multicolumn{1}{l}{\bf Test log-perplexity}
	&\multicolumn{1}{l}{\bf ROUGE-L}
\\ \hline \\
	\textit{cheating} & combined & 1.72975  & 59.3 \\
	\textit{tf-idf} & combined & 2.46645 & 34.2 \\
	\textit{tf-idf} & citations-only  & 3.04299  & 22.6 \\
	\textit{tf-idf} & search-only & 3.56593  & 2.8 \\
	\textit{identity} & combined & 4.80215  & 4.0 \\
\end{tabular}
\end{center}
\end{table}

\textbf{Input Corpus:}
From table \ref{ext-input-perf} we also observe that, unsurprisingly,
the \textit{combined} dataset performs best, but the gaps between it and using
only one of \textit{citations} or \textit{search results} are both significant
and their contributions are
complementary. In subsequent experiments, we report only the \textit{combined} results.



% Results pulled from
% https://docs.google.com/document/d/1wTI2zzkr39KOclO8h8sFAKtxQeqiAwe0QGaXRvottVc/edit
\begin{table}[t]
\caption{Performance of best models of each model architecture
    using the combined corpus and tf-idf extractor. }
\label{abs-perf}
\begin{center}
\begin{tabular}{lll}
\multicolumn{1}{c}{\bf Model} &\multicolumn{1}{c}{\bf Test perplexity} &\multicolumn{1}{c}{\bf ROUGE-L}
\\ \hline \\
\textit{seq2seq-attention, $L=500$}          &  5.04952 & 12.7 \\
\textit{Transformer-ED, $L=500$}          & 2.46645  & 34.2 \\
\textit{Transformer-D, $L=4000$}          & 2.22216  & 33.6 \\
\textit{Transformer-DMCA, no MoE-layer, $L=11000$}          & 2.05159  & 36.2 \\
\textit{Transformer-DMCA, MoE-128, $L=11000$}          & 1.92871  & 37.9 \\
\textit{Transformer-DMCA, MoE-256, $L=7500$}          & 1.90325  & 38.8 \\
\end{tabular}
\end{center}
\end{table}
\begin{figure}[h]
\begin{center}
    \includegraphics[width=0.7\textwidth]{perplexityVsL.png}
\caption{Shows perplexity versus $L$ for tf-idf extraction on combined corpus
	for different model architectures. For T-DMCA, $E$ denotes
	the size of the mixture-of-experts layer.}
\label{perpVsL}
\end{center}
\end{figure}
\textbf{Abstractive model architecture and input length:}
As we see from Table \ref{abs-perf}, \textit{seq2seq-attention} as a baseline
does quite poorly on this task compared to the Transformer architectures.
As seen in Figure \ref{perpVsL}, we observe that the Transformer encoder-decoder, T-ED, architecture
consistently improves in performance until a best of around $L=500-1000$ and is unable to learn at $L=2000$.
This motivated the Transformer-Decoder, which we found could learn and improve up to $L=4000$, before
running out of memory on our machines equipped with 16GB of GPU RAM (NVIDIA P100).
By using the T-DMCA modifications, we were able to train up to $L=11000$ and continued to see
improvements in performance.
We also found the MoE-layer helped performance by adding model capacity at high $L$, for example
dropping log-perplexity from 2.05 to 1.93 at $L=11000$ with 128 experts. Our best model
attempted uses 256 experts at $L=7500$
(we were unable to use 256 experts with $L=11000$ due to memory constraints)
and achieves a perplexity of 1.90,


\textbf{Human Evaluation - Linguistic quality}
We conducted a DUC-style human evaluation of linguistic quality{\footnote{\url{http://duc.nist.gov/duc2007/quality-questions.txt}} of samples from 
    a baseline abstractive (seq2seq), the best extractive (\textit{tf-idf}), and our best T-DMCA models.
    Five different dimensions are assessed: grammaticality, non-redundancy, referential clarity, focus, and structure/coherence. As seen in Table \ref{ling-quality}, the T-DMCA model does statistically significantly better 
    on all dimensions, 
    except on non-redundancy where \textit{tf-idf} does about as well. Overall,
    we observed high fluency and coherence from our best abstractive model.
    Occasionally we observed some repetition of phrases which hurt
    the non-redundancy and structure, but it was much rarer compared with
    the other abstractive method, \textit{seq2seq}. The biggest weakness
    of the extractive method compared with our best abstractive model
    was the lack of structure and coherence in the summaries.
    

\begin{table}[t]
    \caption{Linguistic quality human evaluation scores
    (scale 1-5, higher is better). A score significantly different
(according to the Welch Two Sample t-test, with $p=0.001$)
than the \textit{T-DMCA} model is denoted by *.}
\label{ling-quality}
\begin{center}
\begin{tabular}{llllll}
\multicolumn{1}{l}{\bf Model}
    &\multicolumn{1}{l}{\bf Focus}
	&\multicolumn{1}{l}{\bf Grammar}
    &\multicolumn{1}{l}{\bf \thead{Non-\\redundancy}}
    &\multicolumn{1}{l}{\bf \thead{Referential \\ clarity}}
    &\multicolumn{1}{l}{\bf \thead{Structure and \\ Coherence}}
\\ \hline \\
\textit{T-DMCA (best)} & 4.5 & 4.6 & 4.2 & 4.5 & 4.2 \\
\textit{tf-idf}-only  & 3.0* & 3.6* & 3.9 & 3.2*  & 2.7*  \\
\textit{seq2seq-attention}  & 3.0*  & 3.4*  & 2.1*  & 3.4*  & 2.3* 
\end{tabular}
\end{center}
\end{table}

\textbf{Human Evaluation - side-by-side preference}
We validated our chosen metrics correlate with human preference by conducting
two side-by-side human evaluation experiments, comparing models with large gaps in
perplexity/ROUGE.  We observe in Table \ref{human_eval_table} that human judgment correlates with 
our automatic metrics, but it becomes more difficult to distinguish at the higher-end of model performance.
Details of the human evaluation experimental designs can be found in Appendix \ref{human_eval}.

\begin{table}[t]
\caption{Side-by-side for two models pair with large automatic metric gaps }
\label{human_eval_table}
\begin{center}
\begin{tabular}{lllll}
\multicolumn{1}{l}{\bf Model A} &\multicolumn{1}{l}{\bf Model B}
                                &\multicolumn{1}{l}{\bf ROUGE-L A}
                                &\multicolumn{1}{l}{\bf ROUGE-L B}
                                &\multicolumn{1}{l}{\bf $\frac{\textrm{\#  prefer B}}{\textrm{\# prefer A}}$}
\\ \hline \\
T-ED, $L=100$ & T-ED, $L=500$ &  30.9 & 34.2  & 4.25 \\
T-ED, $L=500$ & T-DMCA-MoE-256, $L=7500$ & 34.2 & 38.8  & 1.5 \\
\end{tabular}
\end{center}
\end{table}

To summarize the quantitative results, we believe the highest impact future work
will be from improving the extractive stage
and extending the decoder-only architectures to learn from larger $L$ while maintaining sufficient
model capacity.

\textbf{Comparison with \cite{Sauper2009}:}
A direct comparison with \cite{Sauper2009} is difficult for three reasons: (a) they report
results only for two small subsets of Wikipedia, Diseases and American Actors; 
 (b) we report on lead generation instead of full-articles;
 (c) we were unable to obtain the exact articles they used as input and output (in particular they make no claim of Wiki-clone detection). However, we
 make a best-effort comparison by finding the subset of articles
 of our test set
 that correspond to Diseases and American Actors, the two categories
 reported on by Sauper \& Barzilay and reporting our ROUGE-1 scores
 (Table \ref{sauper_table}). We observe that we perform better
 on American Actors than Diseases, probably because of the 
 prevalence of the former (and biographies) in Wikipedia compared to the latter in our training set for our single, global model, 
 whereas Sauper \& Barzilay likely benefit from the category-specific templates. On average our ROUGE-1 scores
 are higher but do worse on the less common and 
 somewhat specific disease category.
\begin{table}[t]
    \caption{Comparison of results with \cite{Sauper2009}. Note our results are reported
    for lead section, whereas Sauper \& Barzilay report for articles.}
\label{sauper_table}
\begin{center}
\begin{tabular}{llll}
\multicolumn{1}{l}{}
    &\multicolumn{1}{l}{\bf ROUGE-1 R}
    &\multicolumn{1}{l}{\bf ROUGE-1 P}
    &\multicolumn{1}{l}{\bf ROUGE-1 F1}
\\ \hline
\\
\bf All Wikipedia \\ 
\hline\\
\textit{T-DMCA} (Ours)  & 46  & 53  & 43 \\
\\
\bf Diseases \\
\hline \\
\textit{T-DMCA} (Ours), $n=161$ & 25  & 48 & 29 \\
Sauper \& Barzilay  & 36 & 39 & 37 \\
\\
\bf American Actors \\
\hline \\
\textit{T-DMCA} (Ours), $n=1322$   & 52 & 72 & 54 \\
Sauper \& Barzilay  & 46  & 40 & 41 \\ \\
\end{tabular}
\end{center}
\end{table}


\subsection{Qualitative Discussion}
\begin{figure}[h]
\begin{center}
\fbox{%
    \includegraphics[width=1.0\textwidth]{compareModels.png}
}%
\caption{Shows predictions for the same example from different models. Example model input can
be found in the Appendix \ref{dewey_input}}
\label{compareModels}
\end{center}
\end{figure}
In Figure \ref{compareModels}, we show the predictions from three
different models (using \textit{tf-idf} extraction,
and the \textit{combined} corpus) along with the Wikipedia ground truth. As the perplexity 
decreases we see improvements in the model outputs, in terms of fluency, factual accuracy,
and narrative complexity.  
In particular, the T-DMCA model offers 
a respectable alternative to the Wikipedia version and is more succinct, while mentioning
key facts, such as where the law firm was located, when and how it was formed, and the
rise and fall of the firm.


In manual inspection of model outputs, we noticed an unexpected side-effect: models learn to
translate names from English into multiple languages, e.g. Rohit Viswanath
into Hindi (see Figure \ref{translation}). Although we did not do a systematic evaluation
of the translations, we found they are often correct, and often they are not found in the
Wikipedia article itself. We also verified that in general the translation
is not merely copied from the source, such as example cases where the target language
is the incorrect one (e.g. translation of an English name into Ukrainian).
\begin{figure}[h]
    \centering
\fbox{%
    \includegraphics[width=0.9\textwidth]{translations.png}
}%
\caption{Translation examples from the Transformer-ED, $L=500$.}
\label{translation}
\end{figure}



\subsection{Generating full-Wikipedia articles}
Given that we have shown it is possible to learn sequence transduction models
on combined input-output sequence lengths of approximately 12000 using the
T-D architecture, we show that it is possible to
train a model to generate entire Wikipedia articles.
As a preliminary result, we trained two T-DMCA models: One is trained
to use $L=6000$ reference tokens to predict at most $2192$ article tokens
(longer examples are ignored) and another is conditioned only on
the title and generates articles up to $4000$ tokens long.

We show samples from both models in Appendix~\ref{full_wiki}.
Although the generated articles are not as good as the real Wikipedia
or our lead section samples,
the models can be seen to organize the article into plausible sections
and exhibit global coherence over multi-paragraph text.
The model with access to reference documents inserts factual information
in the generated article. Although we did not focus or tune on
the full-article task we see this as an interesting future work for
abstractive summarization.

